\section{Discussion}
\label{sec:discussion}

\subsection{Why the effect appears on one split and not the other}

The most useful insight this experiment produced is a clean instance of a common
failure mode. Had we reported validation loss and the generalization gap, the
two quantities our hypotheses were literally stated on, we would have
concluded that four register tokens regularize a data-starved ViT-Tiny, with a
consistent direction in every seed and $p = 0.024$. The test split demonstrates
otherwise, and it does so unambiguously: the same arm's test-loss difference is
$+0.002$ with $p = 0.823$.

Three properties of the setup explain the dissociation. The validation split
holds 1{,}000 images, one tenth of the test split, and its between-seed
dispersion is nine times larger. The same split is used to choose the
checkpoint, so the reported validation loss is a minimum over 50 epochs of a
noisy sequence, and anything that changes the shape of that sequence changes its
minimum without necessarily changing the model's quality. And the training loss
does not move at all across arms, which removes the mechanism a genuine regularizer
would have to act through: there is no trade of training fit for validation fit
here, only a different draw from the same distribution. Furthermore, scaling the
training budget to 300 images per class confirms that this dissociation is not
simply an artifact of extreme starvation: when overall performance scales by
$+6.14\%$, the held-out null result persists intact.

We draw a methodological conclusion from this rather than only a negative
empirical one. In small-data ablations, a generalization-gap argument built on a
validation split that also performs model selection is not evidence of
generalization. The gap should be reported, as we do, but the claim has to be
settled on held-out data that played no role in checkpoint selection.

\subsection{Why registers may not transfer to this regime}

Our result does not contradict \cite{darcet2024registers}; it bounds where that
mechanism applies. Four differences separate the regimes.

\textbf{Pretraining mismatch.} Our backbone is ImageNet-pretrained and its
registers are not. The pretrained weights encode 12 blocks of routing behaviour
learned in a sequence that had no register positions in it, and the model is
then given 50 epochs over 9{,}000 images to discover that a new set of tokens is
available as scratch space. In \cite{darcet2024registers} the registers are
present throughout pretraining, over datasets larger by three to four orders of
magnitude. It would be surprising if a routing convention learned over hundreds
of millions of images reorganised itself in this budget, and our entropy
measurements indicate it did not: the layer-wise profile of every register arm
is the pretrained profile.

\textbf{The artifacts are present but unmoved.} Our control arm shows a
patch-norm outlier rate of $1.23\%$, about $2.4$ of 196 patches per image. That
is roughly nine times the $0.13\%$ a three-sigma threshold would select from
Gaussian norms, so the heavy-tailed structure the artifact metric is designed to
detect is genuinely there. What is absent is any response to the treatment: the
rate is $1.25\%$ at $K=8$, and the paired differences are positive for every
register arm. The intervention that was introduced to absorb this structure does
not measurably absorb it here, which is a stronger negative statement than
``there was nothing to remove''.

\textbf{Capacity.} At $d = 192$ with 3 heads, eight registers occupy $4\%$ of a
205-token sequence and compete for the same softmax mass as 196 patches. The
$K=8$ degradation we measure is small but is the study's most statistically
consistent effect on held-out data, and it is the direction H2 predicts.

\textbf{Sequence-level dilution, not parameter count.} It is worth being precise
about the mechanism H2 names. Registers add 1{,}536 parameters at $K=8$; no
plausible capacity argument runs through that number. The dilution, if it is
real, is in the attention distribution rather than in the weights: every query
row must now allocate probability across $K$ additional keys that carry no image
content, and at $d = 192$ the model has little margin to spare.

\subsection{What a positive result would have required}

If registers regularize, the signature should be a training loss that rises
while held-out loss falls. We observe neither. A future test with a real chance
of detecting the effect would pretrain with the registers in place, or fine-tune
long enough for the routing to reorganise, or start from a regime where the
artifact rate is measurably elevated in the first place. Absent one of those
changes, adding register tokens to a pretrained small ViT and fine-tuning on
scarce data is, on this evidence, a no-op at best.

\subsection{An exploratory check: registers without training}

The negative result above concerns registers introduced via fine-tuning. A
separate mechanism, proposed independently of our study, avoids training
entirely: Jiang et al.~\cite{jiang2025} identify the specific neurons
responsible for high-norm outlier tokens and redirect their activity into an
untrained token slot at inference time, with no gradient updates. We ran a
small, exploratory check of this idea against our own $K=0$ checkpoints, using
a simplified single-block reimplementation rather than a full port of their
method.

Three conditions were tested across the same three $K=0$ checkpoints (seeds
42, 1337, 3407): a passive empty token slot with no redirection ($+0.07$ pp
over baseline), active redirection of four identified neurons at block~1
($+0.064$ pp), and at block~9 ($+0.06$ pp). A fourth condition, redirection at
block~11 (the final block), returned results numerically identical to the
passive-slot condition in every seed. This is not an empirical null finding
but an architectural necessity: block~11 is followed by no further
self-attention, and the classification head reads only the [CLS] token, so
any intervention confined to non-CLS positions at the final block has no path
back to the token the model is evaluated on. We report this as a
methodological note rather than evidence about the register mechanism at
block~11.

Among the two valid targets (blocks~1 and~9), active redirection did not
exceed the passive empty-slot condition, and all differences are far smaller
than the between-seed dispersion of the $K=0$ baseline itself ($0.30$ pp). We
treat this as a small additional data point consistent with our main
conclusion rather than independent confirmatory evidence: this check used a
single block, four neurons, and three seeds, with no paired statistical
treatment matching Table~II, and should not be read with the same confidence
as the primary ablation.

\begin{table}[h]
\centering
\caption{Exploratory test-time register check (no training), $K=0$
checkpoints, three seeds. Block~11 is architecturally incapable of affecting
the CLS-read classification head; see text.}
\label{tab:test_time_registers}
\begin{tabular}{lc}
\toprule
Condition & Test Top-1 (\%) \\
\midrule
Baseline ($K=0$, untouched)         & $75.18 \pm 0.30$ \\
Passive empty slot (no redirection) & $75.25 \pm 0.24$ \\
Redirection at block 1              & $75.25 \pm 0.26$ \\
Redirection at block 9              & $75.24 \pm 0.25$ \\
\bottomrule
\end{tabular}
\end{table}
